% Part: incompleteness
% Chapter: theories-computability
% Section: consis-with-q

\documentclass[../../../include/open-logic-section]{subfiles}

\begin{document}

\olfileid[fr]{inc}{tcp}{con}

\olsection{Les théories cohérentes avec $\Th{Q}$ sont indécidables}

Le théorème suivant affirme davantage que la seule
indécidabilité de~$\Th{Q}$ : toute théorie qui ne
contredit pas~$\Th{Q}$ est indécidable.

\begin{thm}
Soit $\Th{T}$ une théorie dans le langage de
l'arithmétique, cohérente avec~$\Th{Q}$,
c'est-à-dire telle que $\Th{T}\cup\Th{Q}$ soit
cohérente. Alors $\Th{T}$ est indécidable.
\end{thm}

\begin{proof}
Rappelons que $\Th{Q}$ possède huit axiomes
$!Q_1$, \dots, $!Q_8$. Nous pouvons les remplacer
par l'axiome unique $!E=!Q_1\land\dots\land !Q_8$.

Supposons $\Th{T}$ décidable et cohérente avec
$\Th{Q}$. Posons, pour les !!{sentence}s arithmétiques,
\[
C=\Setabs{!A}{\Th{T}\Proves !E\lif !A}.
\]
Le test des codes d'!!{sentence}s est calculable. Comme $!E$ est fixé
et que $\Th{T}$ est décidable, $C$ est calculable; les autres codes sont exclus de~$C$.
Montrons qu'il sépare $\Th{Q}$ de $\Th{\bar Q}$,
ce qui contredira leur inséparabilité effective.

Si $!A$ appartient à~$\Th{Q}$, il est démontrable
à partir de ses axiomes. Le théorème de la déduction
donne alors une !!{derivation}, en logique du premier
ordre, de $!E\lif !A$. Cette implication est donc
un théorème de~$\Th{T}$, et $!A\in C$.

Si, en revanche, $!A$ appartient à~$\Th{\bar Q}$,
la logique du premier ordre démontre
$!E\lif\lnot !A$. Si $\Th{T}$ démontrait aussi
$!E\lif !A$, elle démontrerait $\lnot !E$;
$\Th{T}\cup\Th{Q}$ serait alors incohérente.
Notre hypothèse l'exclut. Ainsi $\Th{T}$ ne
démontre pas $!E\lif !A$, donc $!A\notin C$.

Nous avons $\Th{Q}\subseteq C$ et
$\Th{\bar Q}\subseteq\Complement{C}$.
L'ensemble calculable~$C$ serait un séparateur,
contradiction.
\end{proof}

Ce théorème a une portée considérable. Il implique
par exemple le résultat suivant.

\begin{cor}
La logique du premier ordre pour le langage de
l'arithmétique, c'est-à-dire l'ensemble
$\Setabs{!A}{\text{$!A$ est démontrable en logique du premier ordre}}$,
est indécidable.
\end{cor}

\begin{proof}
Cet ensemble est la théorie des conséquences de
$\emptyset$, qui est cohérente avec~$\Th{Q}$.
\end{proof}

\end{document}
